\documentclass[8pt,a4paper,landscape]{extarticle}
\usepackage{multicol}
\usepackage[margin=0.3in]{geometry}
\usepackage{listings}
\usepackage{xcolor}
\usepackage{amsmath, amssymb}
\usepackage{enumitem}

\setlength{\columnsep}{0.2in}
\setlength{\columnseprule}{0.5pt}
\setlength{\parindent}{0pt}
\setlength{\parskip}{2pt}

\lstdefinestyle{python}{
    language=Python,
    basicstyle=\ttfamily\scriptsize,
    keywordstyle=\color{blue}\bfseries,
    stringstyle=\color{teal},
    commentstyle=\color{gray}\itshape,
    showstringspaces=false,
    breaklines=true,
    breakatwhitespace=true,
    tabsize=2,
    aboveskip=2pt,
    belowskip=2pt,
    xleftmargin=0pt
}
\lstset{style=python}

\begin{document}
\begin{multicols*}{3}
\raggedright

\begin{center}
    \Large\textbf{Team Reference Document}
\end{center}

\section*{1. Environment \& Debugging}
\textbf{Current computational power}: $10^8$ ops/sec. Generate all possible answers in background if large code size limit.

\textbf{Python Errors:}
\begin{itemize}[leftmargin=*,noitemsep,topsep=0pt]
\item \texttt{RTE:} array size too small or recursion max depth.
\item \texttt{EOFError, IndexError, SyntaxError}
\item \texttt{KeyError:} access key not in dict.
\end{itemize}

\textbf{Compare File:}
\begin{lstlisting}[language=bash]
python a.py < in.txt > out.txt && fc out.txt o_out.txt
\end{lstlisting}

\textbf{Debug Wrapper:}
\begin{lstlisting}
def debug(f):
  def w(*a, **kw):
    print(f"In: {f.__name__} {a} {kw}")
    res = f(*a, **kw)
    print(f"Out: {f.__name__} = {res}")
    return res
  return w
\end{lstlisting}

% ---------- DATA STRUCTURES ----------
\section*{2. Data Structures}
\subsection*{Union-Find (Kruskal/Frogger)}
\begin{lstlisting}
class UF:
  def __init__(self, n):
    self.parents = [*range(n)]
    self.rank = [0]*n
  def find(self, x):
    if self.parents[x] != x:
      self.parents[x] = self.find(self.parents[x])
    return self.parents[x]
  def union(self, x, y):
    dx, dy = self.find(x), self.find(y)
    if dx == dy: return False
    if self.rank[dx] > self.rank[dy]:
      self.parents[dy] = dx
    elif self.rank[dx] < self.rank[dy]:
      self.parents[dx] = dy
    else:
      self.parents[dx] = dy
      self.rank[dy] += 1
    return True
\end{lstlisting}

\subsection*{Deque}
\begin{lstlisting}
from collections import deque
dq = deque([1, 2, 3])
dq.popleft(); dq.append(4)
\end{lstlisting}

\subsection*{Fenwick Tree / BIT}
\begin{lstlisting}
def getsum(bit, i):
  s = 0; i += 1
  while i > 0:
    s += bit[i]
    i -= i & (-i)
  return s
def updatebit(bit, n, i, v):
  i += 1
  while i <= n:
    bit[i] += v
    i += i & (-i)
def construct(arr):
  n = len(arr)
  bit = [0]*(n+1)
  for i in range(n):
    updatebit(bit, n, i, arr[i])
  return bit
\end{lstlisting}

% ---------- MATH & NUMBER THEORY ----------
\section*{3. Mathematics \& Number Theory}

\subsection*{Prime Sieve ($O(N \log \log N)$)}
\begin{lstlisting}
from math import sqrt
def prime(n):
  k = [True]*(n-2)
  for i in range(2, int(sqrt(n))+1):
    if k[i-2]:
      j = i*i
      while j < n:
        k[j-2] = False; j += i
  return [i+2 for i in range(n-2) if k[i]]
\end{lstlisting}

\subsection*{Fast Powering Algorithm ($O(\log P)$)}
\begin{lstlisting}
def fast(b, p):
  res = 1
  while p:
    if p&1: res *= b
    b *= b
    p >>= 1
  return res
\end{lstlisting}

\subsection*{Fibonacci $O(\log N)$ (Matrix Exponentiation Identity)}
$F_{2k} = F_k(2F_{k+1} - F_k)$, $F_{2k+1} = F_{k+1}^2 + F_k^2$
\begin{lstlisting}
memo = {0:0, 1:1, 2:1}
def fib(n, m=10**9+7):
  if n in memo: return memo[n]
  k = n//2
  if n%2 == 0:
    memo[n] = fib(k)*(2*fib(k+1)-fib(k))%m
  else:
    memo[n] = (fib(k+1)**2 + fib(k)**2)%m
  return memo[n]
\end{lstlisting}

\subsection*{Legendre's Formula}
Number of times a prime $P$ divides $N!$: $s = \sum_{k=1} \lfloor \frac{N}{P^k} \rfloor$
\begin{lstlisting}
def legendre(num, factor):
  s = 0; curr = factor
  while num >= curr:
    s += num // curr
    curr *= factor
  return s
\end{lstlisting}

% ---------- STRING ALGORITHMS ----------
\section*{4. String Algorithms}
\subsection*{KMP All Pattern Matching $O(N+M)$}
\begin{lstlisting}
def kmp(t, p):
  n, m = len(t), len(p)
  l = [0]*m; j = 0
  if m < 1: return [0]*(n>-1)
  if n < 1: return []
  for i in range(1, m):
    while j and p[i]!=p[j]: j=l[j-1]
    j += (p[i]==p[j]); l[i] = j
  i = j = 0; r = []
  while i < n:
    while i<n and j<m and t[i]==p[j]: i,j=i+1,j+1
    if j==m: r.append(i-j); j=l[j-1]
    elif j: j=l[j-1]
    else: i += 1
  return r
\end{lstlisting}

\subsection*{Trie}
\begin{lstlisting}
class Node:
  def __init__(self):
    self.child = dict()
    self.end = False
def build(strings):
  root = Node()
  for s in strings:
    cur = root
    for c in s:
      if c not in cur.child:
        cur.child[c] = Node()
      cur = cur.child[c]
    cur.end = True
  return root
\end{lstlisting}

\subsection*{Regex Guide (\texttt{import re})}
\begin{itemize}[leftmargin=*,noitemsep,topsep=0pt]
  \item \textbf{Syntax}: \verb|*| $0+$, \verb|+| $1+$, \verb|?| $0/1$, \verb|{n,m}| $n \le x \le m$, lazy by \verb|?| (\verb|*?|, \verb|+?|)
  \item \textbf{Classes}: \verb|\d| digits, \verb|\w| words, \verb|\s| whitespace, \verb|\b| word boundaries. Upper is negation.
  \item \textbf{Lookaround}: Positive \verb|(?= )|, Negative \verb|(?! )|. Lookbehind: Pos \verb|(?<= )|, Neg \verb|(?<! )|
  \item \textbf{Flags}: \texttt{re.I} ignorecase, \texttt{re.M} multiline, \texttt{re.X} verbose
  \item \textbf{Methods}: \texttt{re.findall(string, pat)}, \texttt{re.sub(pat, repl, str)}, \texttt{re.match} (start), \texttt{re.search} (first)
\end{itemize}
\begin{lstlisting}
from re import finditer
[a.start() for a in finditer("k","kakakawww")]
\end{lstlisting}

% ---------- GRAPH THEORY ----------
\section*{5. Graph Theory}
\subsection*{Dijkstra (Min-heap)}
\begin{lstlisting}
import heapq
def dijkstra(s, n, graph):
  # graph: dict/list of iterable (v, weight)
  wmap = [float('inf')]*n; wmap[s] = 0
  h = [(0, s)]
  while h:
    d, u = heapq.heappop(h)
    if d > wmap[u]: continue
    for v, w in graph[u]:
      if w+d < wmap[v]:
        wmap[v] = w+d
        heapq.heappush(h, (w+d, v))
  return wmap
\end{lstlisting}

\subsection*{Kahn's Topological Sort}
\begin{lstlisting}
from collections import deque
def kahn(n, graph, indeg):
  q = deque(c for c in indeg if indeg[c]==0)
  order = []
  while q:
    u = q.popleft()
    order.append(u)
    for v in graph[u]:
      indeg[v] -= 1
      if indeg[v] == 0:
        q.append(v)
  return order if len(order)==n else None
\end{lstlisting}

% ---------- GEOMETRY ----------
\section*{6. Geometry}
\subsection*{Circle Intersection Context \& Middle Area}
Area enclosed by circles of radius $a, b, c$. Middle:
\begin{lstlisting}
from math import sqrt, acos, pi
def calc_mid(a,b,c,a2,b2,c2):
  return a2*acos((c2-b2+a2)/(2*c*a)) + \
         b2*acos((c2-a2+b2)/(2*b*c)) - \
         sqrt((c-a+b)*(c+a-b)*(a+b-c)*(c+a+b))/2
\end{lstlisting}

\subsection*{Closest Pair ($O(N \log N)$ Sweep Equivalent)}
\begin{lstlisting}
from math import dist
def closest_pair(coords):
  coords.sort()
  d = dist(coords[0], coords[1])
  for i in range(2, len(coords)):
    p = i-1; new = coords[i]
    while p >= 0:
      pt = coords[p]
      if abs(new[0]-pt[0]) <= d:
        if abs(new[1]-pt[1]) <= 2*d:
          d = min(d, dist(new, pt))
      else: break
      p -= 1
  return d
\end{lstlisting}

% ---------- PYTHON TIPS & GOLFING ----------
\section*{7. Python Tips \& Golfing}

\textbf{Bitwise Odd Check:}
\begin{lstlisting}
if a & 1: print("odd")
\end{lstlisting}

\textbf{Tilde Golfing:}
\begin{lstlisting}
~x = -x - 1; -~x = x + 1; ~-x = x - 1
\end{lstlisting}

\textbf{Reading all Input (Flat):}
\begin{lstlisting}
a,b,*c = map(int, open(0).read().split())
a = [*open(0)]
a = open(0).read().split()
\end{lstlisting}

\textbf{List Flattening \& Matrix Transpose:}
\begin{lstlisting}
import sum
A = sum([[1,2],[3,4]], []) # Flatten
T = [*zip(*A)]             # Transpose Matrix
\end{lstlisting}

\textbf{Shortest Infinite Loops \& Multi-tests:}
\begin{lstlisting}
for _ in [0]*int(input()): pass
for _ in ' '*int(input()): pass
\end{lstlisting}

\textbf{Splat Operator (\texttt{*args}, \texttt{**kwargs}):}
Unpacks iterables: \texttt{add(*[1, 2, 3])}. Kwargs unpacks dicts: \texttt{fn(**\{"name": "A"\})}. Let arguments accept arbitrary items without strict parity.

\textbf{LRUCache Memoization:}
\begin{lstlisting}
from functools import lru_cache
@lru_cache(maxsize=None)
def dp(i): pass
\end{lstlisting}

\end{multicols*}
\end{document}
